\section{背景知识：保护模式下的段式内存管理}

\subsection{从实模式到保护模式}

在实模式下（8086），CPU 使用 \texttt{segment:offset} 计算物理地址：

$$\text{物理地址} = \text{segment} \times 16 + \text{offset}$$

这种方式没有任何保护——任何代码都能访问任何内存。

保护模式改变了规则：段寄存器不再直接存储段基址，而是存储一个 16 位的
\term{段选择子}（selector）。CPU 用 selector 去一张全局表（GDT）中查找
对应的\term{段描述符}（segment descriptor），描述符包含段的基址、大小、
权限等信息。

地址翻译变成两步：
\begin{enumerate}
  \item \textbf{分段}：逻辑地址 $(\text{selector}:\text{offset})$ → 线性地址 $= \text{base} + \text{offset}$
  \item \textbf{分页}（如果开启）：线性地址 → 物理地址
\end{enumerate}

\begin{whybox}
现代 OS 都用\term{平坦模型}（flat model）：所有段的 base=0, limit=4GiB。
这让分段\textbf{不参与地址转换}（线性地址 == offset），把控制完全交给分页。
但 GDT 不能省略——x86 保护模式硬性要求每个段寄存器都指向一个有效描述符。
\end{whybox}

\subsection{段寄存器的"双面"结构}

每个段寄存器（\reg{CS}, \reg{DS}, \reg{SS}, \reg{ES}, \reg{FS}, \reg{GS}）
实际上包含两部分：

\begin{figure}[H]
  \centering
  % figures/ch03/fig01-selector-hidden-cache.tex — auto-extracted from chapter source
\begin{tikzpicture}
  % 可见部分
  \node[memcell, minimum width=3cm, fill=blue!10] (vis) {selector (16 bit)};
  \node[above=0.1cm of vis, font=\small\bfseries] {可见部分（程序可读写）};

  % Hidden part
  \node[memcell, minimum width=8cm, fill=orange!10, right=0.5cm of vis] (hid) {
    base (32) \quad limit (20) \quad access \quad flags
  };
  \node[above=0.1cm of hid, font=\small\bfseries] {不可见部分（hidden cache，CPU 内部）};

  % 箭头
  \draw[pointer, dashed] (vis.east) -- (hid.west)
    node[midway, above, font=\scriptsize] {lgdt / mov ds 时同步更新};
\end{tikzpicture}

  \caption{段寄存器由可见的 selector 和不可见的 hidden cache 组成}
\end{figure}

当你执行 \texttt{mov ds, ax} 时，CPU 做了两件事：
\begin{enumerate}
  \item 把 \texttt{ax}（selector）写入 \reg{DS} 的可见部分
  \item 用 selector 去 GDT 查表，把找到的描述符\textbf{缓存}到 \reg{DS} 的 hidden part
\end{enumerate}

\textbf{之后的每次内存访问}，CPU 直接用 hidden cache 中的 base/limit/权限，
\textbf{不再查 GDT}。这就是为什么改了 GDT 后必须重新加载段寄存器——
否则 hidden cache 里还是旧的描述符。

\begin{thinkbox}
如果 CPU 每次内存访问都重新查 GDT，性能会怎样？一次内存读取原本只需 1 次
总线事务，现在变成 3 次（读 GDT → 读描述符 → 读目标内存）。
hidden cache 就是为了避免这个开销。
\end{thinkbox}

% ============================================================================

